\documentclass[letterpaper]{article}
\usepackage[top=3.5cm,bottom=4cm,left=2cm,right=2cm]{geometry}
\usepackage{textcomp,amsmath,longtable}
\usepackage{latexsym}
\usepackage[utf8]{inputenc}
\usepackage{graphicx}
\usepackage{fancyhdr}
\usepackage{booktabs}
\pagestyle{fancy}
%headers
%\lhead{\begin{picture}(0,0)\put(0,10){\includegraphics[width=2.5cm]{itm}}\end{picture}}

\chead{\LARGE \textsc{Poster Evaluation Form}\\
\small \textsc{Instituto Tecnológico de Morelia}\\
\textsc{Maestría en Ciencias en Ingeniería Electrónica}
}
%\rhead{
%\textsc{Tecnológico Nacional de México}\\
%\textsc{Instituto Tecnológico de Morelia}\\
%\textsc{Departamento de Ingeniería Electrónica}\\}
\lfoot{Ver 2-Rev-2022}
\rfoot{\begin{picture}(0,0)\put(-40,-25){\includegraphics[width=2.5cm]{itm}}
\end{picture}}
\renewcommand{\headrulewidth}{0pt}
\headsep=45pt
%end headers----
%=-=-=-=-=-=-=-=

% begin questionnaires--
% Set up counters for questions and headings
\newcounter{hnumber}
\setcounter{hnumber}{0}
%\stepcounter{hnumber}
\newcounter{qnumber}
\setcounter{qnumber}{0}
\stepcounter{qnumber}
%=-=-=-=-=-=-=-=
% defining answer macro
\newcommand{\yesno}{{\Large ~$\Box$}}

%=-=-=-=-=-=-=-=%=-=-=-=-=-=-=-=
%% we now define the questions
\newcommand{\question}[1]{
    \hfill \relax \thehnumber.\theqnumber\hfill\hfill &\textsf{#1} &{\small\yesno} & {\small\yesno} &{\small\yesno} &{\small\yesno}&{\small\yesno}\stepcounter{qnumber}\\
}

%=-=-=-=-=-=-=-=%=-=-=-=-=-=-=-=
%% we now define the headings of questions
\newcommand{\heading}[1]{
	\stepcounter{hnumber}
	\thehnumber &\multicolumn{4}{l}{\bf\textsf{#1}} \\ \midrule%
	}

%=-=-=-=-=-=-=-=
\newcommand{\divider}{\hline}

%=-=-=-=-=-=-=-=%=-=-=-=-=-=-=-=
% Document's content
\begin{document}
\centering
Presenter's name: \rule{6cm}{0.5pt}\\

\textbf{Instructions:} Please mark the score given for each evaluation criterion below.

\begin{longtable}{c p{8cm}ccccc}
% Main table's header
\textsc{No.} &\textsc{Description}&\footnotesize \textsc{Poor} & \footnotesize \textsc{Fair} & \footnotesize \textsc{Aver}& \footnotesize \textsc{Good}&\footnotesize \textsc{Excellent}\\
& & 1&2&3&4&5\\
%----------------------
% 1.
\heading{Poster's content (25 points)}
\question{Content's clarity}
\question{Content's quality}
\question{Originality and complexity}
\question{Project's impact}
\question{Main points' support}
%&\emph{Comments on Poster Content:}
\\

&&&&&Pts:&\rule{1.5cm}{0.5pt}\\
%----------------------
%2.
\setcounter{qnumber}{0}
\stepcounter{qnumber}
\heading{Poster's organization (25 points)}
\question{Layout: organized, effective, and interesting}
\question{Visual aids and support}
\question{Important information is easy to grasp}
\question{Topic and main objective clearly identified}
\question{Informative \& clear hypothesis}
%&\emph{Comments on Poster Content:}
\\

&&&&&Pts:&\rule{1.5cm}{0.5pt}\\
%----------------------
%3.
\setcounter{qnumber}{0}
\stepcounter{qnumber}
\heading{Delivery (25 points)}
\question{Professional and confident}
\question{Engaged with audience}
\question{Loud and clear voice}
\question{Appropriate technical language}
\question{Response to reviewer questions}
\\
&&&&&Pts:&\rule{1.5cm}{0.5pt}\\

%----------------------
%4.
\heading{Methodology definition (25 points)}
\setcounter{qnumber}{0}
\stepcounter{qnumber}
\question{Please evaluate the proposed methodology considering that the student must define an experimental setup, measurements, procedures, or analysis. Therefore, the methodology results must be used to support the student’s hypothesis. \textbf{Total of 25 points}}


\\&&&&&Pts:&\rule{1.5cm}{0.5pt}\\
\bottomrule
\end{longtable}

\small \noindent \bf Do you have any suggestions for improvements? (use backside)
\\ 
\vspace{2cm}

\texttt{Please note:
 This form will be given to the presenter after the event to provide feedback}


\end{document}